% Erster wissenschaftlicher Entwurf auf Basis des Projektstands vom 09.08.2026.
% Branch: Abschluss-Arbeit; Commit: 437812f7be59ff6ae21325056d158b497bf56f2b.
% Die Marker [Quelle erforderlich] sind bewusst gesetzt, da im Projekt keine
% literatur.bib und keine verifizierbare wissenschaftliche Literatursammlung vorliegt.

\chapter{Einleitung}

Sprachmodelle können den Zugriff auf umfangreiche Informationsbestände erleichtern, wenn Antworten nicht allein aus dem Modellwissen erzeugt, sondern auf nachweisbare Daten und Dokumente gestützt werden. [Quelle erforderlich] Im Versicherungsbereich reicht dafür eine einzelne Informationsquelle häufig nicht aus. [Quelle erforderlich] Die vorliegende Arbeit untersucht deshalb eine Architektur, die strukturierte Bestandsdaten und unstrukturierte Versicherungsdokumente zusammenführt und ihre Antworten unter mehreren Qualitätsdimensionen prüfbar macht.

\section{Motivation und Problemstellung}

Versicherungsunternehmen verarbeiten Informationen, die sich hinsichtlich Struktur, Herkunft und fachlichem Zweck unterscheiden. Kunden-, Vertrags- und Schadendaten werden typischerweise in operativen Systemen strukturiert geführt, während allgemeine Leistungsbeschreibungen, Ausschlüsse und Obliegenheiten in Versicherungsbedingungen und weiteren Dokumenten enthalten sind. [Quelle erforderlich] Eine fachliche Anfrage kann Informationen aus beiden Bereichen voraussetzen. So lassen sich individuelle Policenmerkmale aus einem Customer-Relationship-Management-System (CRM-System) entnehmen, während die Bedeutung einer Deckung oder die für ein Ereignis geltenden Bedingungen anhand der zugehörigen Dokumente zu prüfen sind.

Ein ausschließlich dokumentenbasiertes System kann individuelle Vertragsdaten nicht verlässlich aus allgemeinen Bedingungen ableiten. Umgekehrt enthalten CRM-Daten zwar konkrete Angaben zu einer Person oder Police, begründen für sich genommen aber keine allgemeine Aussage über Deckungsumfang, Ausschlüsse oder Obliegenheiten. Für Anfragen, die beide Arten von Informationen verbinden, muss ein System daher erkennen, welche Quellen benötigt werden, und die ermittelten Fakten kontrolliert zusammenführen.

Large Language Models (LLMs) schaffen dabei neue Zugänge zu heterogenen Informationsbeständen, erzeugen jedoch nicht zwangsläufig fachlich richtige oder vollständige Antworten. Sie können plausibel formulierte, aber sachlich falsche oder nicht belegte Aussagen hervorbringen. [Quelle erforderlich] Retrieval-Augmented Generation (RAG) begrenzt dieses Problem, indem vor der Antwortgenerierung relevante Inhalte aus einem vorgegebenen Bestand abgerufen werden. Die Antwortqualität bleibt jedoch von der Auswahl geeigneter Dokumente, der Relevanz der gefundenen Passagen und ihrer korrekten Verarbeitung durch das Sprachmodell abhängig. [Quelle erforderlich]

Auch vorhandene Quellenangaben genügen nicht automatisch als Qualitätsnachweis. Ein Beleg kann fehlen, auf eine unpassende Passage verweisen oder nur einen Teil der zugehörigen Aussage stützen. Umgekehrt kann eine Antwort ausschließlich belegte Aussagen enthalten und dennoch fachlich unvollständig sein. Für die Bewertung müssen deshalb die fachliche Korrektheit, die Vollständigkeit, die Quellenstützung der Aussagen und die Qualität der Zitationen voneinander unterschieden werden. [Quelle erforderlich]

Neben diesen Qualitätsrisiken entstehen Sicherheitsanforderungen. Missbräuchliche Eingaben können darauf zielen, Systemanweisungen zu umgehen, interne Informationen offenzulegen oder unzulässig auf personenbezogene Daten zuzugreifen. [Quelle erforderlich] Ein für Versicherungsdaten vorgesehenes System muss solche Anfragen erkennen, Daten angemessen schützen und bei unsicheren Verarbeitungsschritten kontrolliert reagieren. [Quelle erforderlich] Zugleich darf eine restriktive Sicherheitslogik legitime fachliche Anfragen nicht pauschal blockieren.

Die technische Umsetzbarkeit allein beantwortet noch nicht, ob eine Architektur praktisch einsetzbar ist. Hierzu gehören ein verlässliches Routing, begrenzte Zugriffsrechte, nachvollziehbare Systementscheidungen, ein kontrolliertes Fehlerverhalten sowie vertretbare Laufzeiten. Diese Anforderungen stehen teilweise in einem Spannungsverhältnis: Zusätzliche Retrieval-, Reranking-, Groundedness- und Sicherheitsstufen können die Qualität und Kontrollierbarkeit erhöhen, zugleich aber die Komplexität und Antwortlatenz vergrößern.

Aus dieser Problemstellung ergeben sich vier miteinander verbundene Bewertungsdimensionen: fachliche Korrektheit, quellenbasierte Nachvollziehbarkeit, Sicherheit gegen missbräuchliche Eingaben und praktische Einsetzbarkeit. Die Arbeit richtet sich daher an folgender Forschungsfrage aus: „Wie kann eine agentische RAG-Architektur für den Versicherungsbereich so entworfen und evaluiert werden, dass sie zugleich fachlich korrekt, quellenbasiert nachvollziehbar (grounded/cited), sicher gegen missbräuchliche Eingaben und praktisch einsetzbar ist?“

\section{Zielsetzung der Arbeit}

Ziel der Arbeit ist der Entwurf, die technische Umsetzung und die systematische Evaluation einer agentischen RAG-Architektur für den Versicherungsbereich. Der entwickelte Prototyp ist als Assistenzsystem für interne fachliche Anfragen ausgerichtet. Er soll allgemeine Informationen aus Versicherungsdokumenten abrufen, individuelle Kunden-, Policen- und Schadendaten aus einem CRM-System einbeziehen und beide Informationsarten bei kombinierten Anfragen nachvollziehbar voneinander trennen.

Die Architektur verbindet dafür mehrere Verarbeitungsschritte. Anfragen werden einem Dokumentenpfad, einem CRM-Pfad, einem kombinierten Pfad oder einer gesperrten Verarbeitung zugeordnet. Der Dokumentenzugriff kombiniert eine schlagwortbasierte Suche mit einer Embedding-basierten Suche und ordnet die gefundenen Passagen durch einen Cross-Encoder neu. Die strukturierten Daten werden über eine read-only angebundene EspoCRM-Instanz abgerufen. Auf dieser Grundlage erzeugt ein LLM die Antwort und versieht Dokumentaussagen sowie individuelle Vertragsinformationen mit unterscheidbaren Quellenangaben. Groundedness-Prüfungen, Guardrails und deterministische Sicherheitsregeln kontrollieren weitere Verarbeitungsschritte. Audit- und Diagnosedaten machen Route, verwendete Evidenz, Prüfergebnisse und Laufzeiten nachträglich untersuchbar.

Die Evaluation soll nicht nur die Antwortqualität isoliert betrachten. Entsprechend der Forschungsfrage wird untersucht, in welchem Umfang das System fachlich korrekte und hinreichend vollständige Antworten erzeugt, seine Aussagen auf tatsächlich vorhandene Quellen stützt, missbräuchliche Eingaben angemessen behandelt und unter den gewählten technischen Bedingungen betreibbar ist. Dazu werden Komponentenmessungen, Datensatz-basierte Evaluationen und ausgewählte End-to-End-Szenarien miteinander verbunden. Die Einordnung der Ergebnisse berücksichtigt dabei ausdrücklich, dass gute Werte in einer Dimension Schwächen in einer anderen nicht ausgleichen.

\section{Forschungsfragen}

Die Untersuchung folgt der verbindlich festgelegten Forschungsfrage:

„Wie kann eine agentische RAG-Architektur für den Versicherungsbereich so entworfen und evaluiert werden, dass sie zugleich fachlich korrekt, quellenbasiert nachvollziehbar (grounded/cited), sicher gegen missbräuchliche Eingaben und praktisch einsetzbar ist?“

Die vier darin enthaltenen Qualitätsdimensionen strukturieren das Evaluationsdesign. Fachliche Korrektheit wird anhand der Unterstützung durch relevante Retrieval-Ergebnisse, des Vergleichs mit Referenzantworten und fachlich definierter Szenarioanforderungen untersucht. Quellenbasierte Nachvollziehbarkeit umfasst sowohl die Stützung einzelner Aussagen durch den bereitgestellten Kontext als auch das Vorhandensein und die Angemessenheit von Dokument- und CRM-Citations. Die Sicherheitsbewertung unterscheidet zwischen dem erfolgreichen Blockieren missbräuchlicher Eingaben und dem Zulassen legitimer Anfragen. Praktische Einsetzbarkeit wird über Betriebsfähigkeit, Antwort- und Komponentenlaufzeiten, Zugriffsbeschränkungen, Fehlerverhalten sowie die erfolgreiche Bearbeitung kombinierter Anwendungsszenarien betrachtet. Die konkreten Metriken und ihre Grenzen werden im Evaluationsdesign operationalisiert und bei der Ergebnisinterpretation getrennt ausgewiesen.

\section{Beitrag der Arbeit}

Der Beitrag dieser Arbeit besteht in der prototypischen Umsetzung und empirischen Evaluation einer routing- und werkzeugbasierten CRM--RAG-Architektur für interne Versicherungsanfragen. Das entwickelte Artefakt verbindet strukturierte, ausschließlich lesend abgerufene EspoCRM-Daten mit unstrukturierten Versicherungsdokumenten und integriert dafür eine regelbasierte Auswahl zwischen CRM-only, RAG-only und Combined, hybrides Keyword- und Vektorretrieval, Cross-Encoder-Reranking, eine herkunftserhaltende Zusammenführung der Evidenz sowie Zitations-, Groundedness- und Sicherheitsprüfungen. Der technische Beitrag liegt in der kontrollierten domänenspezifischen Integration dieser Komponenten und ihrer nachvollziehbaren Orchestrierung, nicht in der Entwicklung neuer Einzelalgorithmen. Der evaluative Beitrag ist eine mehrdimensionale Untersuchung, die Retrieval- und Antwortqualität, Groundedness, Zitation, Sicherheit und Laufzeit durch Datensatzanalysen, Komponentenmessungen und End-to-End-Szenarien getrennt bewertet und ihre Wechselwirkungen sichtbar macht. Nicht beansprucht werden die Entwicklung eines neuen Sprachmodells, Retrieval- oder Reranking-Verfahrens, ein produktionsreifes Versicherungssystem, autonome oder rechtsverbindliche Schadenentscheidungen sowie eine Validierung mit realen Kundendaten.

\section{Methodisches Vorgehen}

Ausgangspunkt der Untersuchung ist die Abgrenzung des fachlichen Problems und der vier Qualitätsdimensionen. Eine Literatur- und Grundlagenanalyse ordnet LLMs, RAG, agentische Orchestrierung, Retrieval- und Reranking-Verfahren, Groundedness, Citations sowie Sicherheitsmechanismen ein. [Quelle erforderlich] Darauf aufbauend werden Anforderungen an Datenzugriff, Antworterzeugung, Nachvollziehbarkeit, Sicherheit und Betrieb abgeleitet und in einen Systementwurf überführt.

Die Implementierung wird anschließend auf Komponenten- und Systemebene untersucht. Für die verschiedenen Dimensionen werden unterschiedliche Tests benötigt: Retrieval- und Referenzmetriken für fachliche Qualität, Claim- und Citation-Prüfungen für die Quellenbindung, gelabelte Angriffs- und legitime Anfragen für die Sicherheitsbewertung sowie Laufzeit-, Readiness-, Berechtigungs- und End-to-End-Prüfungen für die praktische Einsetzbarkeit. Quantitative Ergebnisse werden durch eine qualitative Analyse fachlich relevanter Fehlerszenarien ergänzt. Abschließend werden die Grenzen der Messverfahren und des Prototyps diskutiert und die Befunde auf die zentrale Forschungsfrage zurückgeführt.

\section{Aufbau der Arbeit}

Kapitel 2 erläutert den theoretischen Hintergrund und ordnet verwandte Arbeiten zu RAG, agentischen Architekturen, Quellenbindung, Sicherheit und Evaluation ein. Kapitel 3 beschreibt die Forschungsmethodik, leitet die Anforderungen aus der Forschungsfrage ab und entwickelt den Systementwurf. Kapitel 4 dokumentiert den experimentellen Aufbau, die Implementierung, die verwendeten Datenbestände und die Konfiguration der untersuchten Komponenten. Kapitel 5 stellt die Ergebnisse entlang der vier Qualitätsdimensionen dar und analysiert auffällige Fehlerfälle. Kapitel 6 diskutiert die Befunde, die Aussagekraft der Metriken und die Grenzen des Vorgehens im Hinblick auf die Forschungsfrage. Kapitel 7 fasst die Arbeit zusammen und zeigt Ansatzpunkte für die Weiterentwicklung auf. Daran schließen sich das Literaturverzeichnis im IEEE-Stil und der Anhang mit ergänzenden technischen und evaluativen Artefakten an.
